% To be included
\begin{thebibliography}{60}
\bibitem{1} \rufont{Андринов А. А., Хайкинс Э., Теория колебаний}, \rufont{М. — Л.}, (1937).
\bibitem{2} \rufont{Ляпунов А. М., Общая задача об устойчивости движения}, \rufont{М. — Л.}, (1945), \rufont{2 изд.}
\bibitem{3} \rufont{Четаев Н. Г., Устойчивость движения}, \rufont{М. — Л.}, (1955).
\bibitem{4} \rufont{Дубошин Г. Н., Основы теории устойчивости движения}, \rufont{изд-во МГУ}, (1952).
\bibitem{5} \rufont{Труды II-го Всесоюзного Совещания по теории и методам автоматического регулирования}, \rufont{т. I, II, III, АН СССР}, (1955).
\bibitem{6} \rufont{Лурье А. И., Некоторые нелинейные задачи теории автоматического регулирования}, \rufont{М. — Л.}, (1951).
\bibitem{7} \rufont{Летов А. М., Устойчивость нелинейных регулируемых систем}, \rufont{М. — Л.}, (1955).
\bibitem{8} \rufont{Немыцкий В. В., Степанов В. В., Качественная теория дифференциальных уравнений}, \rufont{М. — Л.}, (1941).
\bibitem{9} \rufont{«Основы автоматического регулирования»}, \rufont{сб. под ред. В. В. Солодовникова}, \rufont{М.}, (1954).
\bibitem{10} \rufont{Конторович М. И., Операционное исчисление и нестанционарные явления в электрических цепях}, \rufont{М. — Л.}, (1949).
\bibitem{11} \rufont{Лурье А. И., Операционное исчисление}, \rufont{М. — Л.}, (1950).
\bibitem{12} \rufont{Диткин В. А., Кузенцов П. И., Справочник по операционному исчислению}, \rufont{М. — Л.}, (1951).
\bibitem{13} \rufont{Бутаков Ъ. В., Колебания}, \rufont{М. — Л.}, (1954).
\bibitem{14} \rufont{Лаврентьев М. А., Шабат Ъ. В., Методы теории функций комплексно переменного}, \rufont{М. — Л.}, (1951).
\bibitem{15} \rufont{Гарднер и Бернс, Переходные процессы в линейных системах}, \rufont{М.}, (1950).
\bibitem{16} \rufont{Джеймс Х., Николс Н, Филлипс Р., Теория следящих систем}, \rufont{М.}, (1951).
\bibitem{17} \rufont{Воронов А. А., Элементы теории автоматического регулирования}, \rufont{М.}, (1951).
\bibitem{18} \rufont{Солодовников В. В., Топлеев Ю. И., Крутикова Г. В., Частотных Метод построения переходных процессов}, \rufont{М. — Л.}, (1955).
\bibitem{19} \rufont{Остославский И. В., Калачев Г. С., Продольная устойчивость и управляемость самолета}, \rufont{М.}, (1951).
\bibitem{20} \rufont{Эакс Н. А., Основы экспериментальной аэродинамики}, \rufont{М.}, (1953).
\bibitem{21} \rufont{Сикорд Ч., Вопросы ракетной техники}, № 6, 78–96 (1952).
\bibitem{22} \rufont{Цыпкин Я. Э., Автоматика и телемеханика}, № 2–3, 107–128 (1946).
\bibitem{23} \rufont{Цыпкин Я. Э., Бромберг П. В., Изв. АН СССР}, №12, 1163–1168 (1945).
\bibitem{24} \rufont{Курош А. Г., Курс высшей алгебры}, \rufont{М. — Л.}, (1952).
\bibitem{25} \rufont{«Корректирующие цепи в автоматике»}, \rufont{сб. статей}, \rufont{М.}, (1954).
\bibitem{26} \rufont{Унттекер, Ватсон, Курс современного анализа}, \rufont{М. — Л.}, (1933).
\bibitem{27} \rufont{Фельдбам А. А., Электрические системы автоматического регулирования}, \rufont{М.}, (1954).
\bibitem{28} \rufont{Крылов А. Н., Соч., АН СССР}, т. V.
\bibitem{29} \rufont{Кочин Ы. Е., Розе Н. В., Теоретическая гидромеханика}, \rufont{М. — Л.}, (1938).
\bibitem{30} \rufont{Мур Дж. Р. и др., Следящие и авторегулируемые системы, отличающиеся от обычных систем с обратной связью}, \rufont{Прикл. мех. и Машиностр., сб. пер.}, № 5, (1953).
\bibitem{31} \rufont{Вознесенский И. И., За советское энергооборудование}, \rufont{сб. статей}, (1934).
\bibitem{32} \rufont{Айзерман М. А., Теория автоматического регулирования двигателей}, \rufont{М. — Л.}, (1952).
\bibitem{33} \rufont{Айзерман М. А., Введение в динамику автоматического регулирования двигателей}, \rufont{М.}, (1950).
\bibitem{34} \rufont{Маккол Л. А., Основы теории сервомеханизмов}, \rufont{М.}, (1947).
\bibitem{35} \rufont{Боголюбов Н. Н., О некоторых статических методах в математической физике}, \rufont{АН УССР}, (1945), \rufont{стр.} 7.
\bibitem{36} \rufont{Красовский А. А., Автоматика И Телемеханика}, 9,1, (1948).
\bibitem{37} \rufont{Поспелов Г. С., Труды ВВА им. Жуковского}, \rufont{вып.} 335, (1949).
\bibitem{38} \rufont{Попов Е. П., Динамика систем автоматического регулирования}, \rufont{М. — Л.}, (1954).
\bibitem{39} \rufont{Гольдфарб Л. С., Автоматика и телемеханика}, № 5, 349; № 6, 413, (1948).
\bibitem{40} \rufont{«Автоматическое регулирование»}, \rufont{сб. статей}, \rufont{М.}, (1954).
\bibitem{41} \rufont{Цыпкин Я. Э., Переходные и установившиеся процессы в импульсных цепях}, \rufont{М.}, (1951).
\bibitem{42} \rufont{Колмогоров А. Н., Фомин С. В., Элементы теории функций и функционального анализа}, \rufont{Издро МГУ}, (1954).
\bibitem{43} \rufont{Ведров В. С., Динамическая устойчивость Самолета}, \rufont{М.}, (1938).
\bibitem{44} \rufont{Фельдбаум А. А., Автоматика и телемеханика}, № 4, (1949).
\bibitem{45} \rufont{Степанов В. В., Курс дифференциальных уравнений}, \rufont{М. — Л.}, (1937).
\bibitem{46} \rufont{Стокер Дж., Нелинейные колебания в механических и электрических системах}, \rufont{М.}, (1952).
\bibitem{47} \rufont{Фельдбам А. А., Автоматика и телемеханика}, № 6, (1953).
\bibitem{48} \rufont{Фельдбам А. А., Автоматика и телемеханика}, № 2, (1955).
\bibitem{49} \rufont{Льюис., Нелейных обратных связи}, \rufont{Вопросы ракетной техники}, № 4, (1954).
\bibitem{50} \rufont{Булгаков Ъ. В., Прикладная теория гироскопов}, \rufont{М.} (1954), \rufont{2 изд.}
\bibitem{51} \rufont{Мак-Лахлан Н. В., Теория И приложения функций Матье}, \rufont{М.}, (1953).
\bibitem{52} \rufont{Россер Дж., Ньютон Р. и Гросс Г., Математическая Теория полета неуправляемых ракет}, \rufont{М.}, (1950).
\bibitem{53} \rufont{Ренкин Р. А., Математическая теория движения неуправляемых ракет}, \rufont{М.}, (1951).
\bibitem{54} \rufont{Гантмахер Ф. Р., Левин М. Л., Прикл. матем. и мех.}, \rufont{т.} X, 301–312, (1947).
\bibitem{55} \rufont{Дреник Р., Вопросы ракетной техники}, \rufont{сборник переводов}, \rufont{М.},5, (1952).
\bibitem{56} \rufont{Фельбаум А. А., Автоматика и телемеханика}, IX, \rufont{вып.} 1, 3–19, (1948).
\bibitem{57} \rufont{Уланов Г. М., Автоматика и телемеханика}, IX, \rufont{вып.} 3, 168–175, (1948).
\bibitem{58} \rufont{Красовский А. А., Труды ВВА им. Жуковского}, \rufont{вып.} 281, (1948).
\bibitem{59} \rufont{Лаврентьев М. А., Люстерник Л. А., Курс вариационного исчисления}, \rufont{М. — Л.}, (1938).
\bibitem{60} \rufont{Автоматы Сб. статей. под ред К. Э. Шеннона и Дж. Маккарти. изд. и-л. Москва}, (1956).
\end{thebibliography}